\section{Inert-prime quadratic parity and the final two tensors}
\label{sec:inert-parity}
The preceding parity theorem uses the integral basis $1,\sqrt q$ and
requires irrational coefficient one. The following extension treats the
case in which $2$ is inert in the full ring of integers and the
Frobenius--Perron dimensions have irrational coefficient two. Its only
extra finite step enumerates characters over the four-element field.

\begin{theorem}[Inert-prime parity rigidity]\label{thm:inert-parity}
Let $R$ be a self-dual rank-five fusion ring with
\[
 s_0=1,\qquad s_i=A_i+2\sqrt q\quad(1\le i\le4),
\]
where $q>1$ is squarefree, $q\equiv5\pmod8$, $A_i\ge0$, and
$1+\sum_{i=1}^4A_i^2=16q$. Put
$M=(N_{jj}^{\,i})_{1\le j,i\le4}$ and
$B=\sum_{i=0}^4 N_i^2$. Suppose that
\begin{enumerate}
\item $\operatorname{ord}_{t=0}(\det(tI-B)\bmod2)=2$;
\item $\ker_{\mathbb F_2}M=\mathbb F_2(1,1,1,1)$;
\item for every ring homomorphism
 $\chi:R\otimes\mathbb F_2\longrightarrow\mathbb F_4$ with
 $\chi(1)=1$ and $\sum_{i=0}^4\chi(X_i)^2=0$, all
 $\chi(X_i)$ lie in $\mathbb F_2$, and exactly three of
 $\chi(X_1),\ldots,\chi(X_4)$ equal one;
\item $N_j$ is irreducible for some $j>0$ with $A_j>0$.
\end{enumerate}
Then every complex fusion-category categorification of $R$ is
spherical and admits a pseudo-unitary Galois conjugate.
\end{theorem}
\begin{proof}
Assume $\C$ categorifies $R$. Choose the self-dual lifts in its
canonical sphericalization, write their dimensions $d_0=1,d_1,\ldots,d_4$,
and let $T_i$ be the signed multiplication matrices. By
\eqref{eq:signed}, $T_i\equiv N_i\pmod2$ and $|T_i|\le N_i$ entrywise.
Their dimension character has field $K=\Q(d_1,\ldots,d_4)$ and
$\delta=1+\sum_i d_i^2$.

Write $D=1+\sum_i s_i^2=32q+4(\sum_i A_i)\sqrt q$.
Since some $A_j>0$, $D$ generates $\Q(\sqrt q)$; its quadratic
minimal polynomial reduces to $t^2$ modulo two. Lemma~\ref{lem:spectral}
and hypothesis (1) show $K=\Q(\sqrt q)$. We also need the stronger
divisibility already implicit in that lemma: the FP character of $R$
inflates to a character of the sphericalization with formal codegree
$2D$. The corresponding simple summand of its induced unit has
dimension $(2\delta)/(2D)=\delta/D$ by
\cite[Theorem 2.13]{Ostrik}. Central-simple dimensions are algebraic
integers, so $\delta/D\in\OO_K$. Since $D\in2\OO_K$, this implies
$\delta\in2\OO_K$.

Now $\OO_K=\ZZ[\omega]$, $\omega=(1+\sqrt q)/2$. Because
$q\equiv5\pmod8$, the reduction of $\omega$ satisfies
$\bar\omega^2+\bar\omega+1=0$; thus $\OO_K/2\OO_K=\mathbb F_4$.
The residues $\bar d_i$ constitute a character $\chi$ of $R$ over
$\mathbb F_4$ because $T\equiv N\pmod2$, and they satisfy
$\sum_i\bar d_i^2=\bar\delta=0$. By (3) they all lie in $\mathbb F_2$,
with exactly three nonunit residues equal to one. Writing
$d_i=x_i+y_i\omega$ with integers $x_i,y_i$, this forces $y_i$ even.
Consequently $d_i=a_i+b_i\sqrt q$ for integers $a_i,b_i$.

The two Galois-conjugate dimension characters are distinct orthogonal
eigenvectors of the real symmetric commuting matrices $T_i$. Their
orthogonality reads
\begin{equation}\label{eq:inert-orthogonal}
 1+\sum_{i=1}^4 a_i^2=q\sum_{i=1}^4 b_i^2.
\end{equation}
Not all $b_i$ vanish, since $K$ is quadratic. Let $g$ be the positive
gcd of the $b_i$, and put $c_i=b_i/g$. Comparing irrational parts in
$d_j^2=\sum_iT_{jj}^{\,i}d_i$ gives
$2a_jb_j=\sum_iT_{jj}^{\,i}b_i$. Divide by $g$, then reduce modulo
two: $Mc=0$. The vector $c$ is primitive; by (2), every $c_i$ is odd.
In particular each $b_i\ne0$.

If $g$ were odd, then all four $b_i$ would be odd. Reduction in
$\OO_K/2$ sends $\sqrt q=2\omega-1$ to one, so
$\bar d_i=\bar a_i+1$. Since exactly three nonunit $\bar d_i$ equal
one, exactly one of the $a_i$ is odd. But
\eqref{eq:inert-orthogonal} modulo four yields
$1+\sum a_i^2\equiv q\sum b_i^2\equiv0\pmod4$, and would require
exactly three odd $a_i$. This contradiction makes $g$ even; hence
$|b_i|\ge2$ for every $i$.

The spectral-radius inequality for $|T_i|\le N_i$ gives
$|d_i|,|d_i'|\le s_i$, and therefore
$|a_i|+|b_i|\sqrt q\le A_i+2\sqrt q$.
Thus $|a_i|\le A_i$. Hypothesis $1+\sum A_i^2=16q$, together with
\eqref{eq:inert-orthogonal}, implies equality throughout:
$|a_i|=A_i$ and $|b_i|=2$ for every $i$.

Choose the real embedding for which $|d_j|=s_j$ at the irreducible
$N_j$ in (4). The Perron-vector argument in the proof of
Theorem~\ref{thm:parity}, applied verbatim to $|T_j|\le N_j$,
then gives $|d_i|=s_i$ for all $i$. Changing the selected lift signs
accordingly gives a signed tensor $U$ with $|U_{ik}^{\,\ell}|\le
N_{ik}^{\,\ell}$ and the common positive dimension character $s$.
Every coefficient of $N-U$ is nonnegative, and its dimension-weighted
sum is zero; hence $U=N$. The selected lifts therefore form a
spherical fusion subcategory mapping tensor equivalently to $\C$.
Under the chosen embedding this structure has positive simple
dimensions, proving the claimed pseudo-unitary Galois conjugate.
\end{proof}

The hypotheses are recognizable from the full tensor, without a
census label. The irrational coefficients of $s_j^2$ give
$A_j=\frac12\sum_{i=1}^4N_{jj}^{\,i}$; the proposed $q$ is then
$(1+\sum A_i^2)/16$. The supplement checks every product identity,
squarefreeness, the characteristic polynomial modulo two, the kernel
of $M$, all $4^4$ possible residue tuples, and irreducibility.

\begin{corollary}\label{cor:inert-frontier}
Neither $\F{18}{5073}$ nor $\F{18}{5076}$ admits a categorification
over $\CC$, with no pivotal hypothesis.
\end{corollary}
\begin{proof}
For the two tensors in \eqref{eq:frontier-tensors}, the vectors
$(A_1,A_2,A_3,A_4)$ are respectively
$(33,34,33,35)$ and $(33,33,34,35)$, and $q=285$. In both cases the
mod-two zero multiplicity of $\det(tI-B)$ is two, and
$\ker_{\mathbb F_2}M=\mathbb F_2(1,1,1,1)$.
With $\bar\omega^2+\bar\omega+1=0$, all normalized residue characters
are as follows; tuples display their four nonunit entries:
\[
\begin{array}{c|l|l}
 &\text{all residue characters}&\sum_i\chi(X_i)^2=0\\ \hline
5073 &(0,1,0,1),(0,1,1,0),(1,0,1,1),(1,1,0,0)&(1,0,1,1)\\
5076 &(0,0,1,1),(1,1,0,1),(\bar\omega,\bar\omega^2,1,0),
 (\bar\omega^2,\bar\omega,1,0)&(1,1,0,1)
\end{array}
\]
Theorem~\ref{thm:inert-parity} applies. The independent exact
spherical-character certificates supplied for these tensors then
exclude the resulting spherical categories. Both finite checks are
recomputed from the tensors in the supplement.
\end{proof}
